% ch10_g §10.12 续 (书页 375–378 / p389–p392)
这个结果与 (10.115) 十分相似。$c_{\mathbf{q}}^* c_{\mathbf{q}}$ 的本征态，比方说记作 $|n_{\mathbf{q}}\rangle$，对应于该模中激发了 $n_{\mathbf{q}}$ 个量子。在简单情形下，反铁磁振子（antiferromagnon）的能量将是
\begin{align*}
\hbar\omega_{\mathbf{q}} &= (A_{\mathbf{q}}^2 - B_{\mathbf{q}}^2)^{\frac{1}{2}} \\
&= \bigl\{ (\beta H_A - 2zSJ)^2 - (2pSJ\langle\cos qa\rangle)^2 \bigr\}^{\frac{1}{2}}, \tag{10.135}
\end{align*}
其中 $\langle\cos qa\rangle$ 是这个量对 $p$ 个最近邻取的平均值，这些近邻与原点相距 $a$，我们假定它们通过交换积分 $J$ 相互作用。如果写成
\begin{equation*}
\hbar\omega_A = \beta H_A, \quad \hbar\omega_J = -2pSJ \tag{10.136}
\end{equation*}
（记得 $J$ 在反铁磁性中是负的），则
\begin{equation*}
\omega_{\mathbf{q}} = \bigl\{ (\omega_A + \omega_J)^2 - (\omega_J\langle\cos qa\rangle)^2 \bigr\}^{\frac{1}{2}} \tag{10.137}
\end{equation*}
就是自旋波谱。

这看上去相当简单。例如，设 $\omega_A \ll \omega_J$，那么
\begin{align*}
\omega_{\mathbf{q}} &\approx \omega_J\bigl\{ 1 - \langle\cos qa\rangle^2 \bigr\}^{\frac{1}{2}} \\
&\approx \frac{1}{\sqrt{3}}\,\omega_J\, .\, qa, \tag{10.138}
\end{align*}
对应于频率对波数呈线性依赖，正如晶格谱 (2.21) 中那样。磁振子对比热的贡献应当按 $T^3$ 变化，也正如 (2.57) 中那样。

让我们试着遵循 (10.122)，计算每个子晶格（sublattice）的磁化强度随温度的变化。平均自旋偏离显然是
\begin{align*}
\langle a_l^* a_l\rangle &= \frac{1}{N}\sum_{\mathbf{q}} \langle b_{\mathbf{q}}^* b_{\mathbf{q}}\rangle \\
&= \frac{1}{N}\sum_{\mathbf{q}} \bigl\langle \cosh^2\theta_{\mathbf{q}}\, c_{\mathbf{q}}^* c_{\mathbf{q}} + \sinh^2\theta_{\mathbf{q}}\, c_{\mathbf{q}} c_{\mathbf{q}}^* \bigr\rangle \\
&= \frac{1}{N}\sum_{\mathbf{q}} \Bigl\{ \bigl(n_{\mathbf{q}} + \tfrac{1}{2}\bigr)\frac{A_{\mathbf{q}}}{\hbar\omega_{\mathbf{q}}} - \tfrac{1}{2} \Bigr\}, \tag{10.139}
\end{align*}
这里用到了变换 (10.127) 与 (10.130)。在此式中 $n_{\mathbf{q}}$ 是第 $q$ 个模中激发的磁振子数；我们自然要用 Bose–Einstein 分布 (10.120) 把它作为温度的函数计算出来。

但是即使在绝对零度，平均自旋偏离也不会消失。在 (10.139) 中取 $n_{\mathbf{q}} = 0$，我们仍然有
\begin{align*}
\langle a_l^* a_l\rangle_{T=0} &= \frac{1}{2N}\sum_{\mathbf{q}} \left( \frac{A_{\mathbf{q}}}{\hbar\omega_{\mathbf{q}}} - 1 \right) \\
&\approx \frac{1}{2N}\sum_{\mathbf{q}} \Bigl\{ \frac{1}{\bigl(1 - \langle\cos qa\rangle^2\bigr)^{\frac{1}{2}}} - 1 \Bigr\} \tag{10.140}
\end{align*}
（在 (10.138) 所描述的简单情形下）。这个和是可以计算的；它的行为如同积分
\begin{equation*}
\int \frac{1}{q}\,\mathrm{d}\mathbf{q} \tag{10.141}
\end{equation*}
只有当 $\mathrm{d}\mathbf{q}$ 是在二维或三维流形上取值时，它才收敛。于是我们证实：一维反铁磁体即使在 $T = 0$ 时也不具有稳定的有序态。这超出了 (10.96) 的结果——后者表明一维 Ising 模型对热涨落是不稳定的；在反铁磁情形，与自旋交换相联系的涨落会摧毁体系所假想的有序基态。

%FIG{197}: 在反铁磁有序态 (a) 中，交换算符 $\mathsf{S}_l^+ \mathsf{S}_{l'}^-$ 产生两个自旋偏离 (b)。

即使在二维和三维情形，每个子晶格的磁矩也达不到其假定的极大值。事实上，由 (10.125) 中的 $\sigma_l$ 所定义的“交替”态并非真正是哈密顿量的本征态。我们可以试着用 (10.106) 的方法来核查这一点；正如在那里指出的，$\mathsf{S}_l^+ \mathsf{S}_{l'}^-$ 项倾向于把一个自旋偏离从一个格点传递到相邻的下一个格点。而在反铁磁情形，它会在两个相邻格点上同时产生自旋偏离——于是基态必定包含由这种组态及类似组态出发所能得到的一切状态的混合。

三维体系的反铁磁基态并非完全有序（当 $S = \tfrac{1}{2}$ 时，每个子晶格大概只能达到其极大磁矩的 90\%），这一事实也反映在能量的期望值上。从 (10.126)、(10.129) 与 (10.134) 可以看到，它是
\begin{align*}
\langle\mathcal{H}\rangle + \mathcal{E}_0 &= \sum_{\mathbf{q}} \bigl( n_{\mathbf{q}} + \tfrac{1}{2} \bigr) \hbar\omega_{\mathbf{q}} - A_{\mathbf{q}} + \mathcal{E}_0 \\
&= -2NpS(S+1)J + \sum_{\mathbf{q}} \bigl( n_{\mathbf{q}} + \tfrac{1}{2} \bigr) \hbar\omega_{\mathbf{q}}, \tag{10.142}
\end{align*}
即便在 $T = 0$、$n_{\mathbf{q}} = 0$ 时，它也大于人们对“Ising”态算出的对角值。磁振子各模中的“零点能”必须包括进去；它使基态的能量升高，基态能量位于 $-2NpJS(S+1)$ 与 $-2NpJS^2$ 之间的某处。

%FIG{198}: 在没有各向异性场时，反铁磁阵列 (a) 的自旋可以在能量不变的情况下转动 (b)。

关于 (10.139) 还有进一步的一点值得注意。按照 (10.138)，$\omega_{\mathbf{q}}$ 在 $q = 0$ 处为零，这样，在平均自旋偏离的计算中这一项的贡献就会是无穷大。我们需要各向异性场（anisotropy field）$H_A$ 来使得
\begin{equation*}
\omega_{\mathbf{q}} \to \sqrt{\bigl( 2\omega_A\omega_J \bigr)} \quad \text{当}\; q \to 0 \tag{10.143}
\end{equation*}
从而避免这个奇点。事实是：我们的体系在其总自旋为零的基态下，对于如图 198 所示的所有自旋的整体均匀转动是不稳定的。各向异性场正是用来稳定量子化方向的。这个场可以非常小——没有 $H_A$，(10.141) 中的极限在 $q = 0$ 处也可以是有限的——但原则上各向异性场必须存在，而且实际上它也确实存在。

实际上，我们可以把频率 (10.143) 作为体系的反铁磁共振（antiferromagnetic resonance）频率观测到。即使 $H_A$ 只有几百高斯，这也是一个很高的频率，但它可以通过施加一个强的外磁场而被带入共振。

反铁磁有序态的磁化率规则不难核查；当然，关于磁振子–声子相互作用等等也已有各种计算。然而，最有意思的一点或许是：如果我们用 (10.139) 去计算有限温度下子晶格的极化，就会发现一个行为如下的表达式
\begin{align*}
\sum_{\mathbf{q}} \frac{\bar{n}_{\mathbf{q}} A_{\mathbf{q}}}{\hbar\omega_{\mathbf{q}}} &\approx \sum_{\mathbf{q}} \frac{kTA_{\mathbf{q}}}{(\hbar\omega_{\mathbf{q}})^2} \\
&\propto \int \frac{1}{q^2}\,\mathrm{d}\mathbf{q} \tag{10.144}
\end{align*}
它在二维情形下于 $q = 0$ 处对数发散。

这样，交换效应在受到热激发时，就足以摧毁二维 Ising 模型的有序反铁磁态；Onsager 结果在这个情形下对真实自旋并不成立（尽管它对铁磁性确实成立）。积分 (10.144) 在三维情形确实收敛，因此反铁磁态在三维下很可能是稳定的——总会有一个由 (10.35) 那样的偶极–偶极项产生的微小各向异性场。但这一点尚未得到严格证明；其正确性的最好明证，也许当属 §2.6 中提到过的中子衍射实验——它们在物理上直接显示出了两个磁子晶格。
